\intro{
\label{first_page}

\textbf{Justification of the choice of the research topic.}
In the National report on the state and prospects of the development of education in Ukraine, informatisation, digitalisation, and intellectualisation of all spheres of life are indicated as one of the basic transformational processes, and the digitalisation of education is a determining factor in its development \cite{UNESCO2023}. Digital history in science is an interdisciplinary approach to the study of history by digital means \cite{CohenRosenzweig2006}. In education and teaching methods, digital history acts as a new approach to the organisation of educational material and the learning process: creating interactive maps, processing large datasets, recognising handwritten texts, and creating infographics all facilitate mastery of educational content, thereby expanding students' cognitive abilities \cite{Wright-Maley2018, Block2021, Tilak2023}.

Elements of digital history are already used in the educational system of Ukraine at all levels. Since 2019, with the onset of the pandemic and later the full-scale Russian invasion, there has been an active move towards online courses and distance education, as well as localised use of various digital tools. In this context, gamification -- the application of game design elements to non-game contexts \cite{Deterding2011} -- has emerged as the most extensively researched digital approach to history education: a systematic review of 74 studies \citep{KorniienkoSemerikov2024} confirmed that gamification dominates the literature on digital tools in history pedagogy \cite{encyclopedia3040089}. Consequently, this dissertation positions gamification as a key subset of the broader digital history methodology, consistent with the UDC code [004.9+930]::378.147.

Digital history and its gamification component address several interconnected challenges:
\begin{enumerate}
\item Formation and development of digital competencies of educational process participants, aligned with the European DigComp~2.0 framework \citep{Vuorikari2016, Titova2025cte};
\item Improvement of the learning process through interactive, evidence-based assessment tools \cite{Sailer2017, Qian2016};
\item Strengthening ties with the European educational research community and advancing integration processes.
\end{enumerate}

Even though the use of Information and Communication Technologies (ICT) in education is thoroughly represented in the studies of the last decade \cite{DichevDicheva2017, KoivistoHamari2019}, the methods of ICT application cannot be directly transferred to history education without adaptation. In the framework of globalisation and informatisation, there is a need for a new, systemic and specialised approach. Digital history allows combining existing methodological assets with new methods, techniques, and tools that exist at the intersection of education, science, and technology \cite{CohenRosenzweig2006, Wright-Maley2018}. At the same time, for use in both higher and general secondary education, this approach needs experimental verification and scientific substantiation \cite{Hattie2009, Tamim2011}, which should first be realised with future specialists as more mobile and independent learners.

Thus, the relevance of the scientific understanding of the research problem is due to contradictions between:
\begin{itemize}
\item the societal demand for digitalisation of education and the insufficient development of a coherent methodology for teaching digital history, including the absence of theoretically grounded, experimentally verified approaches \cite{UNESCO2023, Block2021, Hamari2014, DichevDicheva2017};
\item the importance of forming digital competencies of future specialists in the conditions of digitalisation and the lack of discipline-specific frameworks linking digital competence development to history education \cite{Vuorikari2016, Titova2025cte};
\item the growing body of gamification research in history education (74 studies) and the critically low methodological quality of existing studies (83.8\% high risk of bias per MRQS) \citep{Korniienko2025mrqs}.
\end{itemize}

\textbf{Connection of work with scientific programmes, plans, topics.} The dissertation work was carried out in accordance with the plans of the scientific research work of the Department of Pedagogy of Kryvyi Rih State Pedagogical University and the topic of the complex scientific study ``Methodology of teaching Digital History to future specialists'' (State registration number 0120U101116 dated 26.02.2020).

\textbf{The purpose of the study} is to theoretically substantiate, develop, and experimentally test a digital history teaching methodology for future professionals, including a methodology model that simultaneously operationalises multiple digital history approach types within a single course.

\textbf{The object of research} is the educational process of training future professionals in digital history methodologies within higher education institutions.

\textbf{The subject of the study} is the pedagogical methodology for teaching digital history to future professionals, including its theoretical foundations, practical implementation strategies, and effectiveness assessment.

\textbf{Primary research question.} How can the teaching methodologies of digital history be designed to effectively prepare students for professional careers in digital historical fields while maintaining academic rigour and developing historical thinking?

\textbf{Secondary research questions:}
\begin{enumerate}
  \item How do current theoretical frameworks in digital humanities pedagogy address the specific needs of professional training in digital history?
  \item What methodological quality characterises existing research on gamification in history education, and what gaps need to be addressed? \cite{Korniienko2025mrqs}
  \item What institutional, technological, and wartime conditions affect equitable access to digital history tools in professional training environments?
  \item What design principles for gamification are most effective in the context of history education? \cite{Deterding2011, KorniienkoSemerikov2024}
\end{enumerate}

\textbf{Research hypothesis.} The implementation of a gamification-enhanced digital history teaching methodology will improve student engagement and learning outcomes in history education \cite{Hamari2014, Sailer2017}, as measured by repeated assessment performance, while the systematic analysis of existing research quality will provide evidence-based guidelines for future implementations \cite{Page2021, Korniienko2025mrqs}.

\textbf{Research methods.} To achieve the stated purpose, the following methods were employed:
\begin{itemize}
\item \textit{Theoretical methods:} analysis and synthesis of scientific literature; systematic review following the PRISMA 2020 guidelines \citep{Page2021}; methodological quality assessment using the Mixed Methods Research Quality Synthesis (MRQS) tool; comparative analysis of gamification definitions and frameworks \cite{Deterding2011, DichevDicheva2017}.
\item \textit{Empirical methods:} pedagogical experiment (ex post facto single-group repeated-measures design); quantitative analysis of Google Forms quiz responses; qualitative content analysis of the course chat log; analysis of gamification platform usage.
\item \textit{Statistical methods:} descriptive statistics (mean, median, standard deviation); score trajectory analysis; attendance pattern analysis.
\item \textit{Validation methods:} LLM-assisted quality assessment validation; cross-referencing of automated and manual MRQS scoring \cite{Korniienko2025mrqs, LandisKoch1977}.
\end{itemize}

\textbf{Scientific novelty of the obtained results.}

\textit{For the first time:}
\begin{itemize}
\item a PRISMA-guided systematic review of digital history pedagogy was conducted (46 Scopus
  records, 6 included), establishing the field as nascent~\citep{Korniienko2025cte}
  (Section~\ref{sec:dh-review});
\item a systematic review of gamification in history education synthesised 74 studies,
  revealing definitional ambiguity (14.9\% explicit definitions), predominantly high risk
  of bias (83.8\% per MRQS), and positive but uncertain effects on engagement (70.3\%)
  and learning (51.4\%; $d\,{=}\,0.28$--$0.61$)~\citep{KorniienkoSemerikov2024}
  (Section~\ref{sec:gamification-review});
\item LLM-assisted MRQS quality assessment was validated, establishing the boundaries of
  automation ($\kappa\,{=}\,-0.04$ to $0.21$ for risk-of-bias; 82--97\% for factual
  data)~\citep{Korniienko2025mrqs} (Section~\ref{sec:mrqs});
\item a seven-month pedagogical experiment under active armed conflict combined all five
  DigComp~2.0 competence areas~\cite{Vuorikari2016}, a four-type digital history taxonomy,
  and a five-layer wartime resilience architecture~\cite{Burde2017, UNESCO2023}
  (Chapters~\ref{ch:methodology}--\ref{ch:results});
\item a digital history methodology model was proposed and empirically instantiated,
  simultaneously operationalising all four digital history approach types within a single
  course~\cite{CohenRosenzweig2006, Siemens2004, Connolly2012} (Chapter~\ref{ch:results});
\item the concept of \textit{reward-bridged learning} (RBL) was proposed as a distinct
  position on the game integration spectrum, where educational micro-tasks are embedded
  within entertainment games as optional reward-driven pathways~\cite{RBL2025},
  extending \citeauthor{Deterding2011}'s~\cite{Deterding2011} taxonomy by reversing
  the directionality of game--education integration (Section~\ref{sec:gamification-review}).
\end{itemize}

\textit{Improved:}
\begin{itemize}
\item the content and methods of history teaching using digital tools, through an integrated
  system combining 156 Google Forms quizzes, 7 modular controls, and HistoryLikeGame
  (62 modules, 13 game types, 9 periods, 10 AR scenes) within a TPACK~\cite{Mishra2006}
  and constructivist~\cite{Vygotsky1978} framework (Chapters~\ref{ch:methodology}--\ref{ch:results});
\item the implementation guidelines for gamification in history education through eight
  DigComp-2.0-aligned guidelines~\citep{Titova2025cte} and a 12-item
  CONSORT/PRISMA-aligned reporting checklist (Section~\ref{sec:mrqs}).
\end{itemize}

\textit{Further developed:}
\begin{itemize}
\item the theoretical foundations for digital history pedagogy, including: a three-domain
  model (History $\times$ Digital Humanities $\times$ Pedagogy); a four-type approach
  taxonomy; and a discipline-specific DigComp interpretation grounding all five competence
  areas in historical epistemology~\cite{KorniienkoSemerikov2024, Wineburg2001, Seixas2013}
  (Chapter~\ref{ch:theory});
\item the iterative DBR methodology through the documented v1$\to$v2 HistoryLikeGame
  trajectory (23$\to$62~modules; 1$\to$13~game types)~\cite{Reeves2003, Korniienko2024hlg}
  (Chapter~\ref{ch:methodology});
\item the five-layer wartime resilience architecture for technology-enhanced education in
  conflict-affected settings, validated across 7~months of active armed
  conflict~\cite{Ertmer1999, Burde2017, UNESCO2023} (Chapters~\ref{ch:methodology}--\ref{ch:results}).
\end{itemize}

\textbf{Personal contribution of the education seeker.} All scientific results presented in this dissertation were obtained personally by the author. In co-authored publications:
\begin{itemize}
\item in the paper co-authored with S.~O.~Semerikov on MRQS quality assessment with LLM validation \citep{Korniienko2025mrqs}, the author personally conducted the systematic review, performed the MRQS assessment, designed and executed the LLM validation experiment, and wrote the manuscript; the co-author provided supervision and methodological guidance;
\item in the paper co-authored with L.~O.~Titova, P.~V.~Zahorodko, M.~V.~Moiseienko, and I.~I.~Donchev on gamification and DigComp~2.0 \citep{Titova2025cte}, the author personally developed the gamification platform, designed the implementation guidelines, and contributed to the theoretical framework and manuscript preparation.
\end{itemize}

\textbf{Practical significance of the obtained results.}
The dissertation yielded five practically transferable outputs directly applicable to
Ukrainian history education and to technology-enhanced education in conflict-affected
settings more broadly.
\begin{itemize}
\item \emph{HistoryLikeGame platform} \cite{Korniienko2024hlg}: a publicly accessible,
  offline-capable interactive platform for Ukrainian History NMT preparation,
  comprising 62 modules across 13 game types (quiz, card solitaire, caf\'{e} simulator,
  A-Frame AR, Three.js 3D, Canvas physics, and others) and 9 historical periods.
  The shared infrastructure ($\approx$88\,KB) is designed for minimal-bandwidth deployment
  under wartime and infrastructure-disrupted conditions.
\item \emph{156 Google Forms quizzes} covering the complete Ukrainian History NMT
  curriculum from ancient times to contemporary Ukraine, organised in 36 thematic
  seminars with creative narrative naming~-- ready for reuse or adaptation in any
  Ukrainian History preparatory course.
\item \emph{Reproducible Python analysis pipeline} (6 modules: chat parsing, assessment
  extraction, score analysis, attendance tracking, gamification cataloguing, figure
  generation) enabling transparent, deterministic processing of educational data from
  group chats, Google Forms responses, and HTML game module catalogues.
\item \emph{Eight implementation guidelines for gamified history education}, grounded in
  DigComp~2.0 \cite{Vuorikari2016, Titova2025cte} and the six evidence-based design
  principles from the gamification systematic review \cite{KorniienkoSemerikov2024}, with
  an accompanying 12-item reporting checklist aligned with CONSORT/PRISMA standards
  \cite{Schulz2010CONSORT, Page2021} for researchers planning gamification studies.
\item \emph{Five-layer wartime resilience architecture}: a directly replicable course
  design model for educational continuity under infrastructure disruption, validated
  empirically across 7 months of active armed conflict \cite{Burde2017, UNESCO2023}.
\end{itemize}
The results of the dissertation have been implemented in the educational process of the
Centre for Pre-University Training and Education of Kryvyi Rih State Pedagogical
University.

\textbf{Approbation of dissertation materials.} The main provisions and results of the dissertation were presented and discussed at:
\begin{itemize}
\item the 12th Workshop on Cloud Technologies in Education (CTE 2025), Kryvyi Rih, Ukraine, 2025;
\item the 5th International Conference on History, Theory and Methodology of Learning (ICHTML 2025), Kryvyi Rih, Ukraine, 2025;
\item the 6th International Workshop on Augmented Reality in Education (AREdu 2023), Kryvyi Rih, Ukraine, 2023.
\end{itemize}

\textbf{Publications.} The main results of the dissertation are published in 6 scientific works (Appendix~\ref{app:C}) \cite{Titova2025cte, Korniienko2025mrqs, Korniienko2025cte, Korniienko2025aredu, KorniienkoSemerikov2024, Korniienko2024hlg}, including: 2 articles in journals included in the list of specialised scientific publications of Ukraine; 2 articles in publications indexed by Scopus; 1 preprint; 1 software platform.

\textbf{Structure and scope of the dissertation.}
The dissertation consists of an introduction, three chapters, conclusions, a bibliography, and appendices. The total volume of the dissertation is \pageref{last_page}~pages, of which the main text occupies \the\numexpr\getpagerefnumber{end_main_text}-\getpagerefnumber{first_page}+1\relax~pages. The work contains \totaltables~tables, \totalfigures~figures. The bibliography includes \arabic{citenum@totc}~references.

.
}
